% !TEX program = xelatex
\documentclass[11pt,a4paper]{article}
\usepackage{../../../00_common/preamble}

\title{2024年度 (AY2024) 同志社大学大学院 生命医科学研究科\\
医工学コース 入学試験問題「数学」解答例}
\author{}
\date{}

\begin{document}
\maketitle
\vspace{-3em}

%======================================================================
\section*{第1問}

\subsection*{前提整理}

$\log(1+x)$ の $x=0$ 近傍の Taylor（Maclaurin）展開を求める.
$\dfrac{d}{dx}\log(1+x)=\dfrac{1}{1+x}$ の等比級数展開を項別積分する.

\subsection*{計算}

\paragraph{Step 1: 導関数を等比級数にする.}
$|x|<1$ で
\[
\frac{1}{1+x}=\sum_{n=0}^{\infty}(-1)^n x^n=1-x+x^2-x^3+\cdots.
\]

\paragraph{Step 2: 項別積分する.}
$\log(1+0)=0$ を初期条件として $0$ から $x$ まで積分すると
\[
\log(1+x)=\sum_{n=0}^{\infty}(-1)^n\frac{x^{n+1}}{n+1}
=\sum_{n=1}^{\infty}\frac{(-1)^{n+1}}{n}x^{n}.
\]
\[
\ans{\
\log(1+x)=x-\frac{x^2}{2}+\frac{x^3}{3}-\frac{x^4}{4}+\cdots
=\sum_{n=1}^{\infty}\frac{(-1)^{n+1}}{n}x^{n}\quad(-1<x\le1)\ }
\]
（収束域の端 $x=1$ では交代級数として収束し,$\log2$ を与える.）

検算:$x=1$ で右辺 $=1-\frac12+\frac13-\cdots=\log2=0.6931\ldots$ となり
$\log(1+1)=\log2$ に一致.$\checkmark$

%======================================================================
\section*{第2問}

\subsection*{前提整理}

$C^2$ 関数 $z=f(x,y)$ のラプラシアン $\Delta z=z_{xx}+z_{yy}$ を,
$x=r\cos\theta,\ y=r\sin\theta$ により極座標で表す.連鎖律を用いる.

\subsection*{計算}

\paragraph{Step 1: 1階の連鎖律.}
$z_r=z_x\cos\theta+z_y\sin\theta$,$z_\theta=-z_x\,r\sin\theta+z_y\,r\cos\theta$.

\paragraph{Step 2: 2階を計算する.}
$z_r$ をさらに $r$ で微分すると
\[
z_{rr}=z_{xx}\cos^2\theta+2z_{xy}\cos\theta\sin\theta+z_{yy}\sin^2\theta.
\]
$z_\theta$ を $\theta$ で微分すると
\[
z_{\theta\theta}=r^2\bigl(z_{xx}\sin^2\theta-2z_{xy}\cos\theta\sin\theta+z_{yy}\cos^2\theta\bigr)
-r\bigl(z_x\cos\theta+z_y\sin\theta\bigr).
\]
最後の項は $-r\,z_r$ である.

\paragraph{Step 3: 組み立てる.}
$z_{rr}+\dfrac1r z_r+\dfrac1{r^2}z_{\theta\theta}$ を作ると,$\pm\frac1r z_r$ が相殺し,
交差項も消えて $\cos^2\theta+\sin^2\theta=1$ により $z_{xx}+z_{yy}$ に等しい.したがって
\[
\ans{\
\Delta z=\frac{\partial^2 z}{\partial r^2}+\frac1r\frac{\partial z}{\partial r}
+\frac1{r^2}\frac{\partial^2 z}{\partial\theta^2}\ }
\]

検算:$z=x^2+y^2=r^2$ で左辺 $\Delta z=4$,
右辺 $z_{rr}+\frac1r z_r=2+\frac1r\cdot2r=4$ で一致.$\checkmark$

%======================================================================
\section*{第3問}

\subsection*{前提整理}

$\displaystyle\int_{-\infty}^{\infty}e^{-(x-a)^2/b}\,dx$（$b>0$）を求める.
平行移動とスケール変換で標準 Gauss 積分に帰着する.

\subsection*{計算}

\paragraph{Step 1: 変数変換する.}
$u=\dfrac{x-a}{\sqrt b}$ とおくと $dx=\sqrt b\,du$,$\dfrac{(x-a)^2}{b}=u^2$ で,
積分範囲は $-\infty<u<\infty$ のまま.よって
\[
\int_{-\infty}^{\infty}e^{-(x-a)^2/b}\,dx=\sqrt b\int_{-\infty}^{\infty}e^{-u^2}\,du.
\]

\paragraph{Step 2: 標準 Gauss 積分を代入する.}
$\displaystyle\int_{-\infty}^{\infty}e^{-u^2}\,du=\sqrt\pi$ より
\[
\ans{\ \int_{-\infty}^{\infty}e^{-(x-a)^2/b}\,dx=\sqrt{\pi b}\ }
\]

検算:結果は平行移動 $a$ に依らず,$b$ について $\sqrt b$ に比例する.
$a=0,b=1$ で $\sqrt\pi$ となり標準形に一致.$\checkmark$

%======================================================================
\section*{第4問}

\subsection*{前提整理}

\[
A=\begin{pmatrix}a&b&0&0\\c&d&0&0\\0&0&p&q\\0&0&r&s\end{pmatrix}
\]
はブロック対角行列 $A=\operatorname{diag}(P,Q)$,
$P=\begin{psmallmatrix}a&b\\c&d\end{psmallmatrix}$,
$Q=\begin{psmallmatrix}p&q\\r&s\end{psmallmatrix}$ である.
正則条件と逆行列を求める.

\subsection*{計算}

\paragraph{Step 1: 正則条件.}
$\det A=\det P\cdot\det Q=(ad-bc)(ps-qr)$.$A$ が正則であるためには
$\det A\neq0$,すなわち
\[
\ans{\ ad-bc\neq0\quad\text{かつ}\quad ps-qr\neq0\ }
\]

\paragraph{Step 2: 逆行列.}
ブロック対角行列の逆行列は各ブロックの逆行列を並べたもの.
$2\times2$ の公式 $\begin{psmallmatrix}a&b\\c&d\end{psmallmatrix}^{-1}
=\dfrac{1}{ad-bc}\begin{psmallmatrix}d&-b\\-c&a\end{psmallmatrix}$ を用いて
\[
\ans{\
A^{-1}=\begin{pmatrix}
\dfrac{d}{ad-bc}&\dfrac{-b}{ad-bc}&0&0\\[6pt]
\dfrac{-c}{ad-bc}&\dfrac{a}{ad-bc}&0&0\\[6pt]
0&0&\dfrac{s}{ps-qr}&\dfrac{-q}{ps-qr}\\[6pt]
0&0&\dfrac{-r}{ps-qr}&\dfrac{p}{ps-qr}
\end{pmatrix}\ }
\]

検算:左上ブロックで
$\dfrac{1}{ad-bc}\begin{psmallmatrix}a&b\\c&d\end{psmallmatrix}\begin{psmallmatrix}d&-b\\-c&a\end{psmallmatrix}
=\dfrac{1}{ad-bc}\begin{psmallmatrix}ad-bc&0\\0&ad-bc\end{psmallmatrix}=I_2$.
右下ブロックも同様で,$AA^{-1}=I_4$.$\checkmark$

\end{document}
